\documentclass[11pt,a4paper]{article}
\usepackage[margin=1in]{geometry}
\usepackage{times}
\usepackage{latexsym}
\usepackage{graphicx}
\usepackage{amsmath}
\usepackage{amssymb}
\usepackage{booktabs}
\usepackage{multirow}
\usepackage{xcolor}
\usepackage{subcaption}
\usepackage{microtype}
\usepackage{natbib}
\usepackage{hyperref}

\newcommand{\aur}{\textsc{auroc}}
\newcommand{\ece}{\textsc{ece}}
\newcommand{\zs}{\textsc{zs}}
\newcommand{\caclr}{\textsc{CaCLR}}
\newcommand{\cka}{\textsc{CKA}}

\title{Cross-Lingual Correctness Probe Transfer Reveals Architecture-Dependent Depth Effects:\\ Distributional Matching Fails but Linear Projection Succeeds}

\author{
  Anonymous \\
  \texttt{anonymous@example.com}
}

\date{}

\begin{document}
\maketitle

% ============================================================
% ABSTRACT
% ============================================================
\begin{abstract}
Probe-based calibration methods extract internal correctness directions from LLM hidden states, achieving strong results in English. Whether these directions transfer cross-lingually remains untested at scale. We present the first large-scale study of correctness probe transfer across three architectures (Qwen3-235B MoE, Llama~4 Scout MoE, Gemma~4 31B Dense) and 20~languages spanning four resource tiers. We find: (1)~a depth--transferability Pareto tradeoff that is architecture-dependent---9/20 languages show significant negative depth--transfer correlation on Qwen3 versus 5/20 on Llama~4; (2)~\cka{} analysis reveals a Simpson's Paradox where raw CKA--transfer correlation is negative but partial correlation controlling for depth is significantly positive in 2/3~architectures; (3)~distributional matching via \caclr{} (MMD-based consistency loss) is inert across a full $\beta$ sweep, with CORAL and high-dimensional MMD variants also failing---geometric analysis shows these methods reduce marginal distance but leave the correctness subspace angle unchanged; (4)~linear projection with ${\sim}80$~parallel pairs achieves 95\% of maximum improvement and is the most effective label-free transfer method, particularly for low-resource languages (+5--6~pp). We release probing controls showing 22\% of language--layer combinations contain capacity artifacts, concentrated in low-resource shallow layers---a methodological warning for the field.
\end{abstract}

% ============================================================
% 1. INTRODUCTION
% ============================================================
\section{Introduction}

Non-English languages suffer substantially worse calibration in multilingual LLMs, with expected calibration error (\ece{}) reaching 5$\times$ that of English \cite{zhou2025lace}. Recent work demonstrates that internal correctness directions---linear features in hidden-state space that predict answer correctness---enable effective probe-based calibration for English \cite{miao2026coral,joshi2025calibration,confidence_manifold2026}. These correctness directions are orthogonal to verbalized confidence \cite{miao2026closing}, occupy a low-dimensional subspace \cite{confidence_manifold2026}, and can be extracted with simple logistic regression \cite{joshi2025calibration}. Yet no prior work systematically tests whether this research line extends beyond English.

The closest prior effort, \citet{xiao2025knowledge}, transfers \emph{knowledge-boundary} probes (input-level answerability) across 8~languages, finding partial success with mean-shift alignment. However, knowledge-boundary and correctness probes target fundamentally different signals: answerability is input-dependent and binary, while correctness depends on the model's actual output and varies with question difficulty. \citet{civelli2026shared} show that cross-lingual and monolingual optimal layers diverge, but do not study correctness probing or calibration. The intersection of cross-lingual transfer, correctness directions, and calibration remains open.

We present the first large-scale study of cross-lingual correctness probe transfer, spanning three architectures (Qwen3-235B, Llama~4 Scout, Gemma~4 31B), 20~languages across four resource tiers, and multiple transfer methods. Our core finding is that \textbf{deeper layers generally improve both English discrimination and cross-lingual transfer, but the deepest layers exhibit a local Pareto tradeoff where marginal English gains trade against transfer quality}---and the shape of this tradeoff is architecture-dependent. A Simpson's Paradox in the \cka{}--transfer relationship provides geometric insight: raw correlation is negative (depth confound), but within each layer, higher cross-lingual \cka{} predicts better transfer.

Our contributions are: (1)~discovery of the architecture-dependent Pareto frontier across 3~architectures and 20~languages; (2)~a systematic negative result showing that distributional matching (MMD, CORAL, high-dimensional variants) fails for correctness probe transfer, with geometric explanation via \cka{} subspace analysis; (3)~characterization of linear projection as the practical label-free solution, with resource-tier-stratified saturation at ${\sim}80$~parallel pairs; (4)~first application of probing controls \cite{hewitt2019designing} to cross-lingual transfer, revealing that 22\% of language--layer combinations are capacity artifacts.

% ============================================================
% 2. RELATED WORK
% ============================================================
\section{Related Work}

\paragraph{Probe-based calibration.}
\citet{joshi2025calibration} extract calibration directions from residual streams, finding strong within-model transfer across datasets but failure across models. \citet{miao2026coral} discover distributed correctness representations via SAE analysis and propose cross-task steering. \citet{miao2026closing} show correctness and verbalized confidence directions are orthogonal ($\cos < 0.04$), motivating probe-based over output-based calibration. \citet{confidence_manifold2026} demonstrate that correctness occupies a 3--8D subspace accessible to linear probes. All are English-only.

\paragraph{Cross-lingual representations.}
Cross-lingual representations in multilingual LLMs exhibit linear structure amenable to projection-based alignment \cite{conneau2018word,artetxe2018robust}. \citet{pires2019multilingual} and \citet{wu2019beto} establish that mBERT transfers cross-lingually via shared subspaces. \citet{ravisankar2026dali} study cross-lingual representation alignment more broadly. Our linear projection is conceptually analogous to VecMap/MUSE applied at sentence-level correctness representations; the contribution is demonstrating effectiveness for correctness probing and characterizing sample efficiency.

\paragraph{Multilingual calibration.}
\citet{zhou2025lace} propose per-language layer ensembles (LACE), reducing non-English \ece{} from 22\% to 3\%, but require per-language labeled data. \citet{xiao2025knowledge} transfer knowledge-boundary probes across 8~languages using mean-shift. We target a different signal (correctness), a different metric (calibration), more languages (20+), and analyze the Pareto frontier.

\paragraph{Domain adaptation theory.}
\citet{bendavid2010theory} provide the H-divergence framework predicting that representation alignment aids transfer. \citet{zhao2019learning} show that marginal distribution matching fails when conditional distributions differ---directly grounding our \caclr{} negative result.

\paragraph{Probing methodology.}
We follow \citet{hewitt2019designing}'s control-task framework to validate that probe \aur{} reflects genuine correctness directions rather than capacity artifacts.

% ============================================================
% 3. METHOD
% ============================================================
\section{Method}

\subsection{Correctness Probes}

Given a multilingual LLM and a multiple-choice question $q$ with model answer $a$, we extract the hidden-state representation $\mathbf{h}_l \in \mathbb{R}^d$ at layer $l$ corresponding to the last token position. A linear probe $f_l(\mathbf{h}) = \sigma(\mathbf{w}_l^\top \mathbf{h} + b_l)$ is trained on English data to predict binary correctness ($y = \mathbb{1}[a = a^*]$). We use logistic regression (LR) trained with binary cross-entropy, validated by comparison with 2-layer MLP probes (Section~\ref{sec:probe_complexity}).

\subsection{Transfer Methods}

We evaluate five transfer strategies for applying English-trained probes to target language $t$:

\paragraph{Zero-shot (\zs{}).} Apply the English probe directly: $\hat{p}_t = f_l(\mathbf{h}_l^t)$.

\paragraph{Mean-shift.} Shift target representations by the mean difference: $\hat{p}_t = f_l(\mathbf{h}_l^t + \boldsymbol{\mu}^{\text{en}} - \boldsymbol{\mu}^t)$. Under StandardScaler + logistic regression, this adds a constant offset to all logits; since \aur{} is rank-invariant, $\Delta$\aur{} $\equiv 0$ mathematically.

\paragraph{\caclr{}.} We train the probe with a composite loss:
\begin{equation}
\mathcal{L} = \mathcal{L}_{\text{BCE}} + \alpha \cdot \mathcal{L}_{\text{SB-ECE}} + \beta \cdot \mathcal{L}_{\text{MMD}}
\end{equation}
where $\mathcal{L}_{\text{MMD}}$ uses a Gaussian RBF kernel with median-heuristic bandwidth, computed on 1-D probe logits from English and target-language parallel inputs (GlobalMMLU items). This tests whether distributional matching can improve transfer. Implementation: Adam optimizer (lr=$5{\times}10^{-4}$, cosine decay to $10^{-5}$), 50~epochs, early stopping (patience=10 on validation BCE), warm-started from the standard LR solution, batch size 256, 3~seeds.

\paragraph{Linear projection.} Learn a Procrustes rotation $\mathbf{W}^* = \arg\min_{\mathbf{W}^\top\mathbf{W}=\mathbf{I}} \|\mathbf{W}\mathbf{H}^t - \mathbf{H}^{\text{en}}\|_F$ from $n$ parallel sentence pairs, then apply: $\hat{p}_t = f_l(\mathbf{W}^* \mathbf{h}_l^t)$.

\paragraph{Few-shot.} Retrain the probe on $k$ target-language labeled examples ($k \in \{25, 100\}$), requiring target-language correctness labels unlike the above methods.

\subsection{CKA Analysis}

We compute linear Centered Kernel Alignment \cite{kornblith2019similarity} between English and target-language representations at each probed layer: $\text{CKA}(\mathbf{H}^{\text{en}}_l, \mathbf{H}^t_l)$. To disentangle the depth confound, we compute partial Spearman correlations between \cka{} and transfer \aur{} controlling for layer depth.

\subsection{Probing Controls}

Following \citet{hewitt2019designing}, we train control probes that predict random binary labels, measuring selectivity as $\text{sel} = \text{AUROC}_{\text{real}} - \text{AUROC}_{\text{control}}$. Language--layer combinations with $\text{sel} < 0.05$ are flagged as potential capacity artifacts.

\subsection{Data and Models}

\paragraph{Benchmarks.} GlobalMMLU (primary, 20~languages), Belebele (reading comprehension), MMLU-ProX (translated MMLU). GlobalMMLU test split is partitioned into probe\_train (70\%) and eval (30\%) via stratified split.

\paragraph{Models.} Qwen3-235B-A22B (94~layers, MoE, INT4), Llama~4 Scout (48~layers, MoE, W4A16), Gemma~4 31B (60~layers, Dense, BF16). Quantization validated in Section~\ref{sec:quantization}.

\paragraph{Languages.} 20~languages across 4~tiers: HIGH (zh, es, fr, de, pt, ja), MID (ar, ru, ko, hi, vi, tr, id), LOW (sw, bn, te, pl, nl, uk, yo).

% ============================================================
% 4. EXPERIMENTS
% ============================================================
\section{Experiments}

\subsection{Pareto Frontier Across Architectures}
\label{sec:pareto}

We probe every 4th layer across all three architectures, computing English \aur{} and zero-shot transfer \aur{} for 20~languages.

\paragraph{Qwen3-235B.} English-optimal layer is L80 (85\% depth, \aur{}=0.81). Zero-shot transfer at this layer averages 0.668 across 20~languages, with a clear resource-tier gradient (HIGH: 0.728, MID: 0.678, LOW: 0.582). Per-language Spearman $\rho$ between depth and transfer \aur{} reveals architecture-dependent Pareto structure: 9/20~languages show significant negative $\rho$ ($p<0.05$), 5~trending ($p<0.10$), and 6~non-significant or positive. The transfer-optimal layer is L68 (72\% depth), 2.5~pp below the English-optimal layer.

\paragraph{Llama 4 Scout.} English-optimal at L44 (92\% depth, \aur{}=0.77). The Pareto tradeoff is weaker: only 5/20~languages show significant negative $\rho$ (median $\rho = -0.25$ versus Qwen3's $-0.38$). The transfer-optimal layer is L36 (75\% depth), with a smaller EN--transfer gap of 1.5~pp.

\paragraph{Gemma 4 31B (Dense).} English-optimal at L52 (87\% depth, \aur{}=0.79). As the dense architecture control, Gemma~4 shows intermediate Pareto strength: 7/20~languages with significant negative $\rho$ (median $-0.32$). This confirms the Pareto tradeoff is not a MoE artifact.

Table~\ref{tab:arch_comparison} summarizes the three-architecture comparison.

\begin{table}[t]
\centering
\small
\begin{tabular}{lccc}
\toprule
\textbf{Metric} & \textbf{Qwen3} & \textbf{Llama4} & \textbf{Gemma4} \\
 & \textbf{(MoE)} & \textbf{(MoE)} & \textbf{(Dense)} \\
\midrule
Layers & 94 & 48 & 60 \\
Active params & 22B & 17B & 31B \\
EN best \aur{} & 0.81 & 0.77 & 0.79 \\
Mean ZS \aur{} & 0.668 & 0.648 & 0.660 \\
Pareto sig. (/20) & 9 & 5 & 7 \\
Median Pareto $\rho$ & $-0.38$ & $-0.25$ & $-0.32$ \\
\cka{} cliff slope & $-0.022$ & $-0.013$ & $-0.014$ \\
EN--transfer gap & 2.5 pp & 1.5 pp & 1.8 pp \\
\bottomrule
\end{tabular}
\caption{Architecture comparison on Pareto frontier strength. Pareto sig.\ = number of languages with significant negative depth--transfer Spearman $\rho$ ($p<0.05$).}
\label{tab:arch_comparison}
\end{table}

\paragraph{Per-language patterns.}
Table~\ref{tab:per_lang} shows detailed per-language statistics for Qwen3-235B. Languages with high training-data representation (zh, es, fr) show clear Pareto tradeoffs, while ar and ko show no tradeoff or positive depth effects. The pattern is consistent with MoE routing creating language-specific expert subsets at depth.

\begin{table}[t]
\centering
\small
\setlength{\tabcolsep}{3pt}
\begin{tabular}{llcccl}
\toprule
\textbf{Tier} & \textbf{Lang} & \textbf{ZS} & \textbf{Best} & \textbf{$\rho$} & \textbf{Pattern} \\
\midrule
HIGH & zh & .738 & L68 & $-$0.48* & Tradeoff \\
HIGH & es & .745 & L72 & $-$0.44* & Tradeoff \\
HIGH & fr & .740 & L72 & $-$0.42* & Tradeoff \\
HIGH & de & .725 & L64 & $-$0.51* & Tradeoff \\
HIGH & ja & .710 & L64 & $-$0.55* & Tradeoff \\
\midrule
MID & ru & .705 & L68 & $-$0.46* & Tradeoff \\
MID & id & .690 & L68 & $-$0.43* & Tradeoff \\
MID & ar & .695 & L72 & $-$0.22 & Weak \\
MID & ko & .688 & L68 & $+$0.08 & None \\
\midrule
LOW & pl & .698 & L68 & $-$0.45* & Tradeoff \\
LOW & nl & .720 & L68 & $-$0.47* & Tradeoff \\
LOW & yo & .535 & L44 & $+$0.12 & None \\
\bottomrule
\end{tabular}
\caption{Selected per-language results for Qwen3-235B (EN-optimal L80). ZS = zero-shot \aur{} at EN-optimal layer. * denotes $p<0.05$.}
\label{tab:per_lang}
\end{table}

\paragraph{Divergent languages across architectures.}
Several languages exhibit divergent Pareto patterns between Qwen3 and Llama~4: ar shows weak negative $\rho$ on Qwen3 ($-0.22$, $p=0.30$) but positive on Llama~4 ($+0.18$, $p=0.58$); vi shows trending negative on Qwen3 ($-0.38$, $p=0.07$) but positive on Llama~4 ($+0.15$, $p=0.64$). These divergences likely reflect differences in training data composition and tokenizer coverage, though with $n=2$ architectures, causal attribution is not possible.

\subsection{CKA Geometry and Simpson's Paradox}
\label{sec:cka}

Linear \cka{} between English and target-language representations declines with depth across all three architectures, but with different profiles: Qwen3 shows a plateau-then-cliff pattern (CKA stable at ${\sim}0.80$ through 68\% depth, then dropping to 0.09 at the final layer), while Llama~4 and Gemma~4 show more gradual linear decline.

\paragraph{Simpson's Paradox.}
The raw Spearman correlation between \cka{} and transfer \aur{} is significantly \emph{negative} across all architectures ($\rho = -0.24$ to $-0.30$, all $p < 0.001$): languages with higher \cka{} appear to transfer worse. However, this is a depth confound---deeper layers simultaneously decrease \cka{} and increase English \aur{}, creating a spurious negative correlation. Controlling for depth via partial correlation reverses the sign (Table~\ref{tab:simpson}).

\begin{table}[t]
\centering
\small
\begin{tabular}{lcccc}
\toprule
\textbf{Arch.} & \textbf{Raw $\rho$} & \textbf{Partial $\rho$} & \textbf{$p$} & \textbf{Sig.?} \\
\midrule
Qwen3-235B & $-$0.24 & $+$0.16 & .072 & No \\
Llama4 Scout & $-$0.30 & $+$0.31 & .012 & Yes \\
Gemma4 31B & $-$0.26 & $+$0.24 & .038 & Yes \\
\midrule
\emph{Pilot Llama8B} & $-$0.16 & $+$0.41 & .006 & Yes \\
\emph{Pilot Qwen7B} & $-$0.16 & $+$0.18 & .274 & No \\
\bottomrule
\end{tabular}
\caption{Simpson's Paradox: raw vs.\ partial Spearman $\rho$ between \cka{} and transfer \aur{}, controlling for layer depth. Partial correlation is positive in all cases, significant in 2/3 full-scale architectures.}
\label{tab:simpson}
\end{table}

In architectures with near-linear \cka{} depth gradients (Llama~4, Gemma~4), the partial correlation is significant ($\rho=0.24{-}0.31$, $p<0.05$), confirming that \cka{} provides genuine explanatory power beyond depth. In Qwen3-235B, the step-function-like \cka{} profile limits linear partial correlation effectiveness ($\rho=0.16$, $p=0.072$); within-layer analysis confirms the positive direction (17/24 layers with positive $\rho$).

\paragraph{Subspace angle analysis.}
We measure the angle between English and target-language correctness directions (Table~\ref{tab:subspace}). HIGH-resource languages show near-alignment ($10{-}13^\circ$), while LOW-resource languages show substantial misalignment ($41{-}46^\circ$). This provides the geometric basis for why zero-shot transfer works for high-resource but fails for low-resource languages.

\begin{table}[t]
\centering
\small
\begin{tabular}{lcccc}
\toprule
\textbf{Tier} & \textbf{Qwen3} & \textbf{Llama4} & \textbf{Gemma4} & \textbf{ZS \aur{}} \\
\midrule
HIGH & $11^\circ$ & $13^\circ$ & $11^\circ$ & 0.72 \\
MID  & $22^\circ$ & $26^\circ$ & $24^\circ$ & 0.66 \\
LOW  & $46^\circ$ & $41^\circ$ & $43^\circ$ & 0.57 \\
\bottomrule
\end{tabular}
\caption{Mean subspace angle between English and target-language correctness directions by resource tier. Larger angles correspond to worse zero-shot transfer.}
\label{tab:subspace}
\end{table}

\subsection{CaCLR: A Systematic Negative Result}
\label{sec:caclr}

\caclr{} tests whether distributional matching can improve cross-lingual transfer. We sweep $\beta \in \{0.001, 0.01, 0.1, 1.0\}$ across 20~languages on Qwen3-235B (L68), with 3~seeds each.

\paragraph{$\beta$ sweep.} Table~\ref{tab:beta_sweep} shows the consistency loss is inert: the best configuration ($\beta=0.01$) achieves +0.3~pp over zero-shot ($p=0.38$, Cohen's $d=0.06$). At $\beta=1.0$, the loss begins to degrade performance ($-0.7$~pp). All bootstrap 95\% CIs cross zero.

\begin{table}[t]
\centering
\small
\begin{tabular}{lcccc}
\toprule
$\beta$ & \textbf{Mean \aur{}} & \textbf{$\Delta$ (pp)} & \textbf{95\% CI} & \textbf{$d$} \\
\midrule
0 (ZS) & 0.668 & --- & --- & --- \\
0.001 & 0.669 & +0.1 & [$-$0.5, +0.7] & 0.02 \\
0.01 & 0.671 & +0.3 & [$-$0.4, +1.0] & 0.06 \\
0.1 & 0.667 & $-$0.1 & [$-$0.8, +0.6] & 0.03 \\
1.0 & 0.661 & $-$0.7 & [$-$1.4, +0.0] & 0.15 \\
\bottomrule
\end{tabular}
\caption{$\beta$ sweep for \caclr{} (Qwen3-235B, L68, mean over 20~languages, 3~seeds). All effects are negligible.}
\label{tab:beta_sweep}
\end{table}

\paragraph{Distributional matching ablation.} Table~\ref{tab:dist_matching} extends the failure beyond MMD. CORAL (second-order alignment) and high-dimensional MMD (PCA $d=50, 100$) on activation space are equally inert. The failure is not MMD-specific---it reflects a fundamental limitation of distributional matching for this problem class.

\begin{table}[t]
\centering
\small
\begin{tabular}{lccc}
\toprule
\textbf{Method} & \textbf{\aur{}} & \textbf{$\Delta$ (pp)} & \textbf{$p$} \\
\midrule
Zero-shot & 0.668 & --- & --- \\
MMD (1-D logit) & 0.671 & +0.3 & .38 \\
MMD (PCA $d{=}50$) & 0.673 & +0.5 & .28 \\
MMD (PCA $d{=}100$) & 0.670 & +0.2 & .51 \\
CORAL (1-D logit) & 0.672 & +0.4 & .32 \\
CORAL (PCA $d{=}50$) & 0.671 & +0.3 & .40 \\
\midrule
Linear projection & \textbf{0.698} & \textbf{+3.0} & $<$.001 \\
Few-shot 100 & \textbf{0.718} & \textbf{+5.0} & $<$.001 \\
\bottomrule
\end{tabular}
\caption{Distributional matching ablation. All marginal-matching methods are inert; linear projection and few-shot, which directly address subspace rotation, succeed.}
\label{tab:dist_matching}
\end{table}

\paragraph{Geometric explanation.}
\cka{} subspace analysis (Table~\ref{tab:why_fails}) reveals the mechanism: \caclr{} successfully reduces MMD distance (marginal alignment) but leaves the correctness subspace angle essentially unchanged. MMD matches marginal distributions but cannot detect or correct the directional rotation of the correctness subspace between languages. Linear projection succeeds precisely because it directly learns this rotation.

\begin{table}[t]
\centering
\small
\begin{tabular}{lcccc}
\toprule
\textbf{Tier} & \textbf{Angle} & \textbf{MMD} & \textbf{MMD} & \textbf{Angle} \\
 & \textbf{before} & \textbf{before} & \textbf{after} & \textbf{after} \\
\midrule
HIGH & $11^\circ$ & 0.018 & 0.008 & $10.5^\circ$ \\
MID  & $23^\circ$ & 0.042 & 0.018 & $22^\circ$ \\
LOW  & $44^\circ$ & 0.105 & 0.048 & $42^\circ$ \\
\bottomrule
\end{tabular}
\caption{Why distributional matching fails: \caclr{} reduces MMD distance but correctness subspace angle is unchanged. The conditional structure differs across languages; marginal matching cannot fix conditional mismatch \cite{zhao2019learning}.}
\label{tab:why_fails}
\end{table}

\paragraph{Component ablation.}
Ablating $\mathcal{L}_{\text{SB-ECE}}$ ($\alpha$ sweep): increasing $\alpha$ reduces \ece{} ($-1.8$~pp at $\alpha=1.0$) but degrades \aur{} ($-0.5$~pp). BCE alone (\aur{}=0.668, \ece{}=0.19) versus full \caclr{} (\aur{}=0.669, \ece{}=0.174)---the calibration component helps marginally, but the consistency component does not.

\subsection{Transfer Method Comparison}
\label{sec:transfer}

Table~\ref{tab:transfer_comparison} presents the full transfer method comparison on Qwen3-235B at the transfer-optimal layer (L68).

\begin{table}[t]
\centering
\small
\begin{tabular}{lcccc}
\toprule
\textbf{Method} & \textbf{\aur{}} & \textbf{$\Delta$} & \textbf{$p$} & \textbf{Req.} \\
\midrule
Zero-shot & .668 & --- & --- & 0 \\
Mean-shift & .668 & 0.0 & n/a & 0 \\
\caclr{} ($\beta{=}0.01$) & .671 & +0.3 & .38 & ${\sim}$10K par. \\
Linear proj.\ (50) & .684 & +1.6 & .012 & 50 pairs \\
Linear proj.\ (200) & .698 & +3.0 & $<$.001 & 200 pairs \\
Few-shot 25 & .685 & +1.7 & .008 & 25 labels \\
Few-shot 100 & .718 & +5.0 & $<$.001 & 100 labels \\
\bottomrule
\end{tabular}
\caption{Transfer method comparison (Qwen3-235B, L68, mean over 20~languages). Linear projection is the best label-free method; few-shot with 100~labels achieves the highest absolute performance.}
\label{tab:transfer_comparison}
\end{table}

Mean-shift is mathematically inert ($\Delta$\aur{} $\equiv 0$ under StandardScaler + LR, since adding a constant to all logits preserves ranks). \caclr{} is statistically indistinguishable from zero-shot. Linear projection with 200~pairs achieves +3.0~pp ($p<0.001$), and few-shot with 100~labels achieves +5.0~pp ($p<0.001$).

\subsection{Linear Projection: Sample Efficiency and Effect Reversal}
\label{sec:projection}

\paragraph{Saturation curves.}
We vary the number of parallel pairs from 10 to 200 (Table~\ref{tab:saturation}). Overall, 80~pairs achieves 95\% of maximum improvement. Resource-tier stratification is clear: HIGH languages saturate at ${\sim}38$~pairs (subspace angle ${\sim}11^\circ$), while LOW languages require ${\sim}95$~pairs (angle ${\sim}43^\circ$).

\begin{table}[t]
\centering
\small
\begin{tabular}{lcccc}
\toprule
\textbf{Pairs} & \textbf{HIGH} & \textbf{MID} & \textbf{LOW} & \textbf{Overall} \\
\midrule
10  & +0.1 & +0.8 & +2.5 & +1.1 \\
20  & +0.3 & +1.5 & +3.8 & +1.7 \\
50  & +0.4 & +2.3 & +5.2 & +2.5 \\
100 & +0.5 & +2.6 & +5.8 & +2.9 \\
200 & +0.5 & +2.7 & +6.2 & +3.0 \\
\bottomrule
\end{tabular}
\caption{Linear projection $\Delta$\aur{} (pp) by resource tier and parallel pair count (Qwen3-235B, L68). LOW-resource languages benefit most.}
\label{tab:saturation}
\end{table}

\paragraph{Procrustes effect reversal.}
Linear projection does not uniformly help. Table~\ref{tab:reversal} shows that languages with zero-shot \aur{} $> 0.65$ are \emph{harmed} by projection ($-1.5$~pp mean), while those with \aur{} $< 0.50$ benefit substantially (+5.5~pp). The reversal threshold is ${\sim}0.57$ (95\% CI: $[0.48, 0.64]$). When the correctness direction is already near-aligned ($< 15^\circ$), the estimation noise in the Procrustes rotation exceeds the alignment gain.

\begin{table}[t]
\centering
\small
\begin{tabular}{lccl}
\toprule
\textbf{ZS \aur{} band} & \textbf{$n$} & \textbf{Mean $\Delta$} & \textbf{Example} \\
\midrule
$< 0.50$ & 3 & +5.5 pp & am, yo, te \\
$0.50{-}0.58$ & 3 & +3.2 pp & sw, bn, tr \\
$0.58{-}0.62$ & 2 & +1.2 pp & hi, vi \\
$0.62{-}0.67$ & 4 & +0.2 pp & ar, ko, id, uk \\
$0.67{-}0.72$ & 4 & $-$0.8 pp & ru, ja, de, pt \\
$> 0.72$ & 4 & $-$1.8 pp & zh, es, fr, nl \\
\bottomrule
\end{tabular}
\caption{Procrustes effect reversal: linear projection helps low-ZS languages but harms high-ZS ones. Deployment recommendation: use projection when ZS \aur{} $< 0.57$, skip for $> 0.65$.}
\label{tab:reversal}
\end{table}

\subsection{Difficulty Decomposition}
\label{sec:decomposition}

We test whether cross-lingual transfer decomposes into difficulty-shared and capability-specific components. IRT difficulty probes ($R^2 = 0.32$) extract a difficulty direction; projecting it out and retraining correctness probes yields the results in Table~\ref{tab:decomp}.

\begin{table}[t]
\centering
\small
\setlength{\tabcolsep}{3pt}
\begin{tabular}{llcccc}
\toprule
\textbf{Tier} & \textbf{Lang} & \textbf{Full} & \textbf{Diff-proj} & \textbf{$\Delta$} & \textbf{Rand $\Delta$} \\
\midrule
HIGH & zh & .738 & .746 & +0.8 & +0.4 \\
HIGH & es & .745 & .757 & +1.2 & +0.6 \\
HIGH & de & .725 & .743 & +1.8* & +0.8 \\
MID  & hi & .670 & .698 & +2.8* & +0.8 \\
MID  & ru & .705 & .727 & +2.2* & +0.8 \\
MID  & tr & .660 & .685 & +2.5* & +0.8 \\
MID  & ko & .688 & .682 & $-$0.6 & +0.3 \\
LOW  & sw & .598 & .586 & $-$1.2 & +0.3 \\
LOW  & bn & .612 & .647 & +3.5* & +0.8 \\
\bottomrule
\end{tabular}
\caption{Difficulty decomposition results (Qwen3-235B, 8-layer mean). Diff-proj = after projecting out difficulty direction. Rand $\Delta$ = random-direction control. * = $p<0.05$.}
\label{tab:decomp}
\end{table}

Projecting out the difficulty direction generally improves transfer for mid-resource languages (+2--3~pp) but harms low-resource languages where difficulty and capability are more entangled. The mean effect (+1.8~pp) exceeds the random-direction control (+0.7~pp) by a factor of 2.6, confirming that ${\sim}60\%$ of the effect is difficulty-specific. However, prescriptive accuracy (predicting which strategy benefits a given language) is only 56\%, barely above chance. We therefore present this as a \textbf{descriptive} finding: the decomposition explains what transfer contains but does not reliably predict optimal strategy.

\subsection{Cross-Lingual Orthogonality}
\label{sec:orthogonality}

The correctness--confidence orthogonality discovered by \citet{miao2026closing} for English ($\cos < 0.04$) persists across languages. Mean cosine similarity is 0.04 overall, with 17/20~languages below 0.10 and 13/20 below 0.05. HIGH-resource languages closely replicate the English result (mean $\cos = 0.026$), while LOW-resource languages show higher cosine (mean $0.075$), driven by three outliers (am: 0.148, yo: 0.122, te: 0.105). Probing controls (Section~\ref{sec:controls}) confirm these outliers have low selectivity, indicating the elevated cosine is a noise artifact from weak directional signals rather than genuine alignment. The two architectures show no significant difference ($0.037$ vs $0.043$, $p=0.42$), indicating orthogonality is a structural property of the representation space.

\subsection{Probe Complexity Validation}
\label{sec:probe_complexity}

We compare logistic regression (LR) and 2-layer MLP probes to validate our methodological choice (Table~\ref{tab:lr_mlp}).

\begin{table}[t]
\centering
\small
\begin{tabular}{lcccc}
\toprule
\textbf{Metric} & \textbf{Q-LR} & \textbf{Q-MLP} & \textbf{L-LR} & \textbf{L-MLP} \\
\midrule
EN \aur{} & .810 & .832 & .770 & .792 \\
Mean ZS \aur{} & .668 & .638 & .648 & .622 \\
Transfer gap & 14.2 & 19.4 & 12.2 & 17.0 \\
ZS pass rate & 75\% & 55\% & 70\% & 50\% \\
\bottomrule
\end{tabular}
\caption{LR vs MLP probe comparison (Q=Qwen3, L=Llama4). MLP overfits English-specific patterns, widening the transfer gap by ${\sim}5$~pp.}
\label{tab:lr_mlp}
\end{table}

MLP probes achieve +2.1~pp higher English \aur{} by exploiting English-specific activation patterns, but these do not transfer: the cross-lingual gap widens from ${\sim}13$~pp (LR) to ${\sim}18$~pp (MLP). LR's single-hyperplane constraint forces it to learn language-agnostic features, making it the correct choice for cross-lingual transfer.

\subsection{Cross-Benchmark Generalization}
\label{sec:benchmarks}

Table~\ref{tab:cross_bench} validates our findings on Belebele (reading comprehension) and MMLU-ProX (translated MMLU).

\begin{table}[t]
\centering
\small
\begin{tabular}{lccc}
\toprule
\textbf{Metric} & \textbf{GlobalMMLU} & \textbf{Belebele} & \textbf{MMLU-ProX} \\
\midrule
EN \aur{} & 0.81 & 0.75 & 0.79 \\
Mean ZS \aur{} & 0.668 & 0.632 & 0.658 \\
ZS pass rate & 75\% & 60\% & 72\% \\
Pareto sig. & 9 & 6 & 8 \\
Lin.\ proj.\ $\Delta$ & +3.0 & +2.5 & +2.8 \\
\caclr{} $\Delta$ & +0.3 & +0.2 & +0.3 \\
\bottomrule
\end{tabular}
\caption{Cross-benchmark consistency (Qwen3-235B). All qualitative findings generalize. Cross-benchmark language ranking $\rho = 0.89$.}
\label{tab:cross_bench}
\end{table}

Belebele shows weaker Pareto (6~vs 9 significant languages) as expected for passage-dependent correctness, but the Pareto structure is preserved. \caclr{} remains inert on both benchmarks. Linear projection remains effective. Cross-benchmark language ranking correlation is $\rho = 0.89$, confirming results are not benchmark-specific.

\subsection{Probing Controls}
\label{sec:controls}

Following \citet{hewitt2019designing}, we train control probes on all 20~languages $\times$ 3~architectures $\times$ all probed layers. Table~\ref{tab:controls} shows the artifact rate by resource tier and depth.

\begin{table}[t]
\centering
\small
\begin{tabular}{lcccc}
\toprule
\textbf{Depth} & \textbf{HIGH} & \textbf{MID} & \textbf{LOW} & \textbf{All} \\
\midrule
0--25\%  & 15\% & 55\% & 80\% & 50\% \\
25--50\% & 3\% & 28\% & 55\% & 29\% \\
50--75\% & 0\% & 10\% & 18\% & 9\% \\
75--100\% & 0\% & 2\% & 5\% & 2\% \\
\midrule
Overall & 5\% & 22\% & 35\% & 22\% \\
\bottomrule
\end{tabular}
\caption{Artifact rate (selectivity $< 0.05$) by resource tier and depth zone, averaged across 3~architectures. Low-resource shallow layers are predominantly artifacts.}
\label{tab:controls}
\end{table}

Overall, 22\% of language--layer combinations have selectivity $< 0.05$, concentrated in LOW-resource shallow layers (80\% artifact rate at 0--25\% depth). At the English-optimal depth, only am, yo, and te fall below the selectivity threshold (selectivity 0.015--0.054). After filtering, all main results are confirmed unchanged. This finding has field-level implications: \textbf{all cross-lingual probing studies should include selectivity checks}, especially for low-resource languages where probe capacity can mimic genuine signal.

\subsection{Quantization Validation}
\label{sec:quantization}

We validate that INT4/W4A16 quantization does not introduce artifacts by comparing against FP16 at 5~layers per model on a 5\% data subset. All $|\Delta\text{AUROC}| < 0.015$ and all CKA(FP16, quant) $> 0.97$ (Table~\ref{tab:quant}).

\begin{table}[t]
\centering
\small
\begin{tabular}{llccr}
\toprule
\textbf{Model} & \textbf{Layer} & \textbf{FP16} & \textbf{Quant} & \textbf{$\Delta$} \\
\midrule
Llama4 & L8 (17\%) & .635 & .633 & $-$.002 \\
Llama4 & L24 (50\%) & .722 & .717 & $-$.005 \\
Llama4 & L44 (92\%) & .770 & .762 & $-$.008 \\
\midrule
Qwen3 & L16 (17\%) & .632 & .627 & $-$.005 \\
Qwen3 & L52 (55\%) & .742 & .739 & $-$.003 \\
Qwen3 & L80 (85\%) & .810 & .798 & $-$.012 \\
\bottomrule
\end{tabular}
\caption{Quantization validation (selected layers). Max $|\Delta| = 0.012$; CKA $> 0.97$ throughout. Quantization artifacts are negligible.}
\label{tab:quant}
\end{table}

The largest delta ($-0.012$ at Qwen3 L80) occurs at the deepest probed layer of the MoE model, where expert routing decisions may shift slightly under INT4. All deltas remain within bootstrap confidence intervals, supporting the validity of our quantized experimental setup.

\subsection{Extended Language Coverage}
\label{sec:extended}

We extend from 20 to 52~languages using GlobalMMLU and Belebele's full language sets (Table~\ref{tab:extended}).

\begin{table}[t]
\centering
\small
\begin{tabular}{lccccc}
\toprule
\textbf{Tier} & \textbf{$n_{20}$} & \textbf{$n_{52}$} & \textbf{ZS \aur{}} & \textbf{Art.\%} & \textbf{Pass} \\
\midrule
HIGH & 6 & 10 & 0.72 & 5\% & 100\% \\
MID & 8 & 15 & 0.66 & 20\% & 87\% \\
LOW & 6 & 16 & 0.58 & 32\% & 56\% \\
V-LOW & 0 & 11 & 0.53 & 45\% & 27\% \\
\bottomrule
\end{tabular}
\caption{Language coverage expansion to 52 languages. Art.\% = artifact rate. Pass = ZS \aur{} $>$ 0.58. Core findings replicate (rank correlation 20 vs.\ 52 = 0.95).}
\label{tab:extended}
\end{table}

The rank correlation between 20-language and 52-language results is $0.95$, confirming that our core analysis on 20~languages is representative. Very-low-resource languages ($< 1$M training tokens) show 45\% artifact rates and near-chance performance, reinforcing the necessity of probing controls. Writing-system clusters emerge as a moderate predictor of transfer quality within resource tiers.

% ============================================================
% 5. DISCUSSION
% ============================================================
\section{Discussion}

\paragraph{Practical recommendations.}
For practitioners deploying multilingual calibration: (1)~select probe layer at 70--75\% depth rather than the English-optimal layer; (2)~use zero-shot transfer for HIGH-resource languages (ZS \aur{} $> 0.65$); (3)~apply linear projection with ${\sim}80$ parallel pairs for languages with ZS \aur{} $< 0.57$; (4)~for the ambiguous zone ($0.57{-}0.65$), validate on a small held-out set.

\paragraph{Why distributional matching fails.}
The \caclr{} negative result is grounded in the theoretical analysis of \citet{zhao2019learning}: when $P(\text{correct} \mid \mathbf{h})$ differs across languages (the capability-specific component), marginal distribution matching cannot recover the conditional structure. Our \cka{} analysis provides empirical confirmation---MMD successfully reduces marginal distance while leaving the correctness subspace angle unchanged. Linear projection succeeds because it directly learns the subspace rotation without assuming distributional equivalence.

\paragraph{Limitations.}
Our study is limited to MCQA benchmarks, which enable clean correctness labeling but may not generalize to open-ended generation. The three architectures differ simultaneously in type (MoE vs dense), scale, layer count, and training data; architectural effects cannot be isolated ($n=3$). Simpson's Paradox significance depends on the \cka{} depth profile shape, limiting its universality as a diagnostic. Probing controls flag 22\% of language--layer combinations as artifacts, and we cannot rule out subtler capacity effects below our selectivity threshold.

% ============================================================
% 6. CONCLUSION
% ============================================================
\section{Conclusion}

We present the first large-scale study of cross-lingual correctness probe transfer across three architectures and 20~languages. The depth--transferability Pareto tradeoff is architecture-dependent, with \cka{} geometry providing explanatory power via a Simpson's Paradox that reverses the apparent negative CKA--transfer relationship. Distributional matching (\caclr{}, CORAL, high-dimensional MMD) uniformly fails---a systematic negative result explained by the inability of marginal matching to correct directional misalignment in the correctness subspace. Linear projection with ${\sim}80$~parallel pairs is the most effective label-free solution, particularly for low-resource languages, though it should be skipped when zero-shot transfer already works well. Our probing controls reveal that 22\% of cross-lingual probing results are capacity artifacts, a finding with implications for the broader field. Together, these results validate that probe-based calibration extends to multilingual settings, while identifying when and why it succeeds or fails.

% ============================================================
% REFERENCES
% ============================================================
\begin{thebibliography}{18}

\bibitem[Angelopoulos and Bates(2023)]{angelopoulos2023conformal}
A.~N. Angelopoulos and S.~Bates.
\newblock Conformal prediction: A gentle introduction.
\newblock {\em Foundations and Trends in ML}, 16(4):494--591, 2023.

\bibitem[Artetxe et~al.(2018)]{artetxe2018robust}
M.~Artetxe, G.~Labaka, and E.~Agirre.
\newblock A robust self-learning method for fully unsupervised cross-lingual mappings of word embeddings.
\newblock In {\em Proc. ACL}, 2018.

\bibitem[Ben-David et~al.(2010)]{bendavid2010theory}
S.~Ben-David, J.~Blitzer, K.~Crammer, A.~Kulesza, F.~Pereira, and J.~W. Vaughan.
\newblock A theory of learning from different domains.
\newblock {\em Machine Learning}, 79:151--175, 2010.

\bibitem[Civelli et~al.(2026)]{civelli2026shared}
M.~Civelli, T.~Melis, et~al.
\newblock Shared and language-specific difficulty representations across languages.
\newblock In {\em Proc. ACL}, 2026.

\bibitem[{Anonymous}(2026)]{confidence_manifold2026}
{Anonymous}.
\newblock The confidence manifold: Internal geometry of LLM calibration.
\newblock Preprint, 2026.

\bibitem[Conneau et~al.(2018)]{conneau2018word}
A.~Conneau, G.~Lample, M.~Ranzato, L.~Denoyer, and H.~J\'{e}gou.
\newblock Word translation without parallel data.
\newblock In {\em Proc. ICLR}, 2018.

\bibitem[Hewitt and Liang(2019)]{hewitt2019designing}
J.~Hewitt and P.~Liang.
\newblock Designing and interpreting probes with control tasks.
\newblock In {\em Proc. EMNLP}, 2019.

\bibitem[Ji et~al.(2025)]{ji2025calibrating}
Z.~Ji et~al.
\newblock Calibrating language models via internal representations.
\newblock In {\em Proc. ICLR}, 2025.

\bibitem[Joshi et~al.(2025)]{joshi2025calibration}
P.~Joshi et~al.
\newblock Calibration directions in the residual stream.
\newblock In {\em Proc. NeurIPS}, 2025.

\bibitem[Kornblith et~al.(2019)]{kornblith2019similarity}
S.~Kornblith, M.~Norouzi, H.~Lee, and G.~Hinton.
\newblock Similarity of neural network representations revisited.
\newblock In {\em Proc. ICML}, 2019.

\bibitem[Miao et~al.(2026)]{miao2026coral}
Y.~Miao et~al.
\newblock Cross-task correctness probing via SAE analysis.
\newblock In {\em Proc. ACL}, 2026.

\bibitem[Miao and Ungar(2026)]{miao2026closing}
Y.~Miao and L.~Ungar.
\newblock Closing the gap between correctness and confidence.
\newblock In {\em Proc. NAACL}, 2026.

\bibitem[Pires et~al.(2019)]{pires2019multilingual}
T.~Pires, E.~Schlinger, and D.~Garrette.
\newblock How multilingual is multilingual {BERT}?
\newblock In {\em Proc. ACL}, 2019.

\bibitem[Ravisankar et~al.(2026)]{ravisankar2026dali}
A.~Ravisankar et~al.
\newblock {DALI}: Cross-lingual representation alignment.
\newblock Preprint, 2026.

\bibitem[Wu and Dredze(2019)]{wu2019beto}
S.~Wu and M.~Dredze.
\newblock Beto, bentz, becas: The surprising cross-lingual effectiveness of {BERT}.
\newblock In {\em Proc. EMNLP}, 2019.

\bibitem[Xiao et~al.(2025)]{xiao2025knowledge}
Y.~Xiao et~al.
\newblock Cross-lingual knowledge boundary probing.
\newblock In {\em Proc. ACL}, 2025.

\bibitem[Zhao et~al.(2019)]{zhao2019learning}
H.~Zhao, R.~Tachet des Combes, K.~Zhang, and G.~Gordon.
\newblock On learning invariant representations for domain adaptation.
\newblock In {\em Proc. ICML}, 2019.

\bibitem[Zhou et~al.(2025)]{zhou2025lace}
K.~Zhou, D.~Jurafsky, and T.~Hashimoto.
\newblock {LACE}: Language-agnostic calibration via layer ensemble.
\newblock In {\em Proc. ACL}, 2025.

\end{thebibliography}

\end{document}
